\section{Boh}

 [...]

\[
    \mathcal{H}_c = \bigoplus_{h, \bar{h}} N_{h, \bar{h}} \mathrm{Vir}_c(h) \otimes \bar{\mathrm{Vir}}_{c}(\bar{h}),
\]

[...]

Let us consider the unitarity of the Virasoro generators. We have
\[
    L_n^\dagger = L_{-n},
\]
it should be an hermiticity condition....

We can assume the normalization of primary tates as
\[
    \bra{h}\ket{h} = 1, \quad \bra{h}\ket{h^{\prime}} = \delta_{h,h^{\prime}}.
\]
If we have this primary state, its scalar product is given by the two-point function as we have seen. what about the product of the descendant states? We can write for the secondary state
\[
    \bra{h} L_{-1}^\dagger L_{-1} \ket{h} = \bra{h} L_1 L_{-1} \ket{h},
\]
and since the commutator is given by
\[
    [L_1,\,L_{-1}] = \dots = 2 L_0,
\]
we can write the previous expression as
\[
    \begin{aligned}
        \bra{h} L_{-1}^\dagger L_{-1} \ket{h} & = \bra{h} (2 L_0 + L_{-1} L_1) \ket{h}                   \\
                                              & = 2 h \bra{h}\ket{h} + \bra{h} L_{-1} L_1 \ket{h} = 2 h.
    \end{aligned}
\]
This is a positive number, so we can say that the state $L_{-1} \ket{h}$ is a physical state (its norm is positive). Thus we have found a first condition for the unitarity of the theory, which is that the conformal dimension \(h \geq 0\) must be positive.

There are many interesting situations of non unitary CFTs or minimal models, in which we will not have this constarint of positivity for \(h\). But we will not discuss this topic here.

We have also that the central charge \(c\) must be positive since \dots

\subsection{Gram Matrix}

The Gram matrix is defined as the matrix of scalar products of the descendant states:
\[
    G_{ij}^{(N)} = \bra{v_i^{(N)}}\ket{v_j^{(N)}},
\]
where \(\ket{v_i^{(N)}}\) are the descendant states at level \(N\).

As we have already seen, the first level is given by the state \(L_{-1} \ket{h}\), and its norm is \(2h\):
\[
    G^{(1)} = \bra{h} L_1 L_{-1} \ket{h} = 2h \geq 0,
\]
while already at the second level we have two descendant states, \(\ket{1} = L_{-2} \ket{h}\) and \(\ket{2} = L_{-1}^2 \ket{h}\), and we can compute
\[
    \begin{aligned}
        \bra{1}\ket{1} & = \bra{h} L_{2}L_{-2} \ket{h} = \bra{h} \left(L_{-2}L_2 + 4L_0 + \frac{c}{2}\right)\ket{h} = 4h + \frac{c}{2}, \\
        \bra{1}\ket{2} & = \bra{h} L_{2}L_{-1}^2 \ket{h} = \dots = 6h,                                                                  \\
        \bra{2}\ket{2} & = \bra{h} L_{-1}^2 L_1^2 \ket{h} = \dots = 4h(2h + 1).
    \end{aligned}
\]
\TODO{Finish the calculation of the Gram matrix.}

Thus if we were to compute the Gram matrix at level 2, we would find
\[
    G^{(2)} = \begin{pmatrix}
        4h + \frac{c}{2} & 6h         \\
        6h               & 4h(2h + 1)
    \end{pmatrix},
\]
and if we were to stop at level three\TODO{do it}, we would find a \(3 \times 3\) matrix, and so on.

From this matrix we can obtain some useful information about the theory. For example, we can compute its determinant, which will give us insight about the linear independence of the descendant states. If the determinant is zero, it means that there is a linear dependence among the descendant states, which can lead to the presence of null states in the theory. Null states are states that have zero norm and are orthogonal to all other states in the theory, including themselves. They play an important role in the structure of CFTs and can lead to constraints on the allowed values of \(h\) and \(c\).

If \(h=0\) we have the identity operator, which corresponds to the vacuum state. And the first descendant state is given by \(L_{-1} \ket{0}\), which is the application of the translation generator on the vacuum. Since the vacuum is invariant under translations, we have \(L_{-1} \ket{0} = 0\), which means that the first descendant state is a null state. This is consistent with the fact that the determinant of the Gram matrix at level 1 is zero when \(h=0\).

In the conformal family of the identity operator (equivalently in the verma module of the vacuum) the first seconday state is absent (null state), while the second descendant state is given by \(L_{-2} \ket{0}\), which is an operator of dimension 2, and it is the energy-momentum tensor \(T(z)\). This is consistent with the fact that the determinant of the Gram matrix at level 2 is zero when \(h=0\) and \(c\) is arbitrary.

Let us indeed compute the determinant of the Gram matrix at level 2:
\[
    \det G^{(2)} = (4h + \frac{c}{2})(4h(2h + 1)) - (6h)^2 = 16 h^3 + 8 h^2 - 2 h c = 2h (16h^2 (2c-10)h + c).
\]
Usually, when solving for the unitarity constraints, we look for the values of \(h\) and \(c\) such that the determinant of the Gram matrix is positive, which ensures that all descendant states have positive norm. If we find a particular set of \(h\) and \(c\) which annihilates the determinant, it means that there is a null state at that level, i.e. a particular linear combination of descendant states at that level that has zero norm. It is orthogonal to all other states in the theory, including itself. If we want to build an irreducible representation of the Virasoro algebra, we need to quotient out by the null states, which means that we need to set the null states to zero in the theory. But a null vector gives birth to an entire submodule of the Verma module, which is generated by the action of the Virasoro generators on the null state, and is made up of null states. Thus, if we want to build an irreducible representation, we need to quotient out by the entire submodule generated by the null state.

One can also find another null state at an higher level, which generates another submodule, which can intersect with the previous one, and so on. This is the structure of the Verma module, which can be reducible or irreducible depending on the presence of null states. It is a difficult problem to find and quotient out all the null states in a given Verma module, but it is an important problem, since it allows us to classify the representations of the Virasoro algebra and to understand the structure of CFTs.

If we know the general form of the Gram Matrix determinant at level \(N\), we can find the values of \(h\) and \(c\) for which the determinant is zero, which gives us the unitarity constraints on the theory. This general form is given by the Kaz determinant formula, which is a formula that gives the determinant of the Gram matrix at level \(N\) in terms of the conformal dimension \(h\) and the central charge \(c\). The Kaz determinant formula is given by
\[
    \det G^{(N)} = \prod_{r,s \geq 1}^{rs\leq N} (h - h_{r,s}(c))^{P(N-rs)}, \quad h_{r,s}(c) = \frac{[(m+1)r - ms]^2-1}{4m(m+1)},
\]
where \(P(N)\) is the number of partitions of the integer \(N\), and \(m\) is a parameter related to the central charge by the formula
\[
    c = 1 - \frac{6}{m(m+1)}.
\]

[\textit{graph h vs c }]

at each fixed value of \(c\) we have a discrete set of values of \(h\) for the same theory, and we want to find the values of \(h\) and \(c\) for which the determinant is positive, which gives us the unitarity constraints on the theory. We can already exclude the region of the parameter space where \(h < 0\) or \(c < 0\), since we have already seen that these values lead to negative norm states. Thus, we can focus on the region of the parameter space where \(h \geq 0\) and \(c \geq 0\).

In this region, we can find the values of \(h\) and \(c\) for which the determinant is zero, which gives us the unitarity constraints on the theory. The Kaz determinant has a polynomial structure, and the zeros of the determinant correspond to the values of \(h\) and \(c\) for which there are null states in the theory. The idea is to find out the roots of the determinant, and there are different ways to do this. The most common way is to parametrize the central charge \(c\) in terms of a parameter \(m\) as we have seen, and then to find the values of \(h_{r,s}(c)\) for which the determinant is zero at each fixed value of \(c\). This gives us a discrete set of values of \(h\) for each fixed value of \(c\), which correspond to the unitarity constraints on the theory.

We get some useful constraints: the two parameters \(r\) and \(s\) are positive integers, and the product \(rs\) must be less than or equal to \(N\), which means that we can only have a finite number of null states at each level. Then they have to satisfy
\[
    1 \leq r \leq m-1, \quad 1 \leq s \leq m,
\]
thus we can write
\[
    (r,\,s) = (m-r,\,m+1-s),
\]
which leads us to a table of values for \(r\) and \(s\) for each fixed value of \(m\), which gives us the unitarity constraints on the theory. This creates different parabulas for each value of \(m\) on the left side of the constant \(c=1\) which intersect the line \(c=1\) in their minimum. Where the parabulas intersects among each other, we have unitary theories with the same central charge but different conformal dimensions.

We know from other results that for \(c>1\) al theories are unitary, while for \(c<1\) we have a discrete set of unitary theories, which are the minimal models. The minimal models are characterized by the values of \(m\) and the corresponding values of \(h_{r,s}(c)\) for which the determinant is zero, which gives us the unitarity constraints on the theory. The minimal models are important because they are exactly solvable CFTs, and they have a rich structure of primary fields and correlation functions. They also have applications in statistical mechanics and condensed matter physics, where they describe critical phenomena in two dimensions.

Systems can have other symmetries apart from the conformal symmetry, and these symmetries can lead to additional constraints on the theory, which can further restrict the allowed values of \(h\) and \(c\) and expand the structure of the algbera (thus the theory). Indeed for \(c>1\) there is an huge amount of theories, which becomes tractable only when we have additional symmetries, which can be used to classify the theories and to find their correlation functions. But for \(c<1\) we have a discrete set of theories, which are the minimal models, and they are already classified by the values of \(m\) and the corresponding values of \(h_{r,s}(c)\) for which the determinant is zero, which gives us the unitarity constraints on the theory (with additional symmetries). Thus, for \(c<1\) we have a complete classification of unitary CFTs, while for \(c>1\) we have an infinite number of theories, which can be classified only with additional symmetries.

With \(m=2\) we have only one unitary theory, which has only the vacuum state and is trivial. From \(m=3\) we have the first non-trivial unitary theory, and the central charges start from \(c=1/2\) and increase as \(m\) increases, approaching \(c=1\) as \(m \to \infty\).

For example we have the Ising model with a distribution of vacancies (in some sites on the lattice we have a vacancy, which means that there is no spin in that site), which has an interesting phase diagram, and it has a tri-critical point.

Let us consider the case \(m=3\), we can compute a table of values for \(r\) and \(s\) at \(c = \frac{1}{2}\)
\[
    [\dots]
\]
and we get only 3 values of \(h\), which correspond to the three primary fields of the Ising model, which are the identity operator, the spin operator, and the energy operator. The identity operator has \(h=0\), the spin operator has \(h=1/16\), and the energy operator has \(h=1/2\). These three primary fields generate the entire spectrum of the theory, and all other fields can be obtained as descendants of these primary fields. The correlation functions of these fields can be computed using the conformal symmetry and the structure of the Verma modules, which allows us to find exact results for the correlation functions in this theory.

We should imagine to have this Kaz table (left table) multiplied by another identical Kaz table for the anti-holomorphic sector (right table), which gives us the full spectrum of the theory, which is given by the tensor product of the holomorphic and anti-holomorphic sectors. Thus, for each value of \(c\) we have a discrete set of values of \(h\) and \(\bar{h}\) for which the determinant is zero, which gives us the unitarity constraints on the theory. And once we get \(h\) and \(\bar{h}\) we can compute the conformal dimension \(\Delta = h + \bar{h}\) and the spin \(s = h - \bar{h}\) of the primary fields, which gives us the full spectrum of the theory. Looking at the universality classes given by these tables, we strongly suspect that this is the Ising model (as well as the free-majorana fermion theory, which is equivalent to the Ising model).

we can have a bosonic description of the Ising model, which is given by a \(\phi^4\) lagrangian, as well as the fermionic description, which is given by a free Majorana fermion lagrangian. These two descriptions are equivalent and dual.